% paper/sections/11c_conservation.tex
%
% §11.2  保存則の検証（Conservation Properties）
%   §11.2.1 運動エネルギー減衰（Taylor-Green 渦）
%   §11.2.2 非圧縮条件の維持
% ═══════════════════════════════════════════════════════════════════════════════
\subsection{保存則の検証}
\label{sec:ns_conservation}
% ═══════════════════════════════════════════════════════════════════════════════

NS ソルバが物理的に正しい保存則を再現することを，
Taylor--Green 渦の減衰問題により検証する．

% ─────────────────────────────────────────────────────────────────────────────
\subsubsection{運動エネルギー減衰}
\label{sec:energy_conservation}
% ─────────────────────────────────────────────────────────────────────────────

%% 実験結果は exp11_3_tgv_energy.py の実行後に挿入

\textbf{設定：}
計算領域 $[0, 2\pi]^2$，周期境界条件，単相流（$\rho = 1$，$\nu = 0.01$）．
初期条件は Taylor--Green 渦：
\begin{align}
  u(x,y,0) &= \sin x\,\cos y, &
  v(x,y,0) &= -\cos x\,\sin y, \\
  p(x,y,0) &= -\tfrac{1}{4}(\cos 2x + \cos 2y).
\end{align}
厳密解は $\bu(t) = \bu(0)\,e^{-2\nu t}$，$p(t) = p(0)\,e^{-4\nu t}$．

\textbf{評価指標：}
\begin{itemize}[noitemsep]
  \item 運動エネルギー $E_k(t) = \frac{1}{2}\int_\Omega |\bu|^2\,d\Omega$ の
        理論減衰率 $E_k(t)/E_k(0) = e^{-4\nu t}$ に対する相対誤差
  \item 非圧縮性 $\|\bnabla\cdot\bu\|_\infty$ の時間推移
\end{itemize}

\textbf{合格基準：}
\begin{itemize}[noitemsep]
  \item 運動エネルギー相対誤差 $< 10^{-4}$（$t = 1$ において）
  \item $\|\bnabla\cdot\bu\|_\infty < 10^{-10}$（全時刻で維持）
\end{itemize}

%% TODO: exp11_3 結果テーブルと図をここに挿入

% ─────────────────────────────────────────────────────────────────────────────
\subsubsection{離散非圧縮条件の定量評価}
\label{sec:div_free}
% ─────────────────────────────────────────────────────────────────────────────

上記 Taylor--Green 渦シミュレーションにおける $\|\bnabla\cdot\bu\|_\infty$ の
時間推移を記録する．CCD-PPE の投影ステップが各タイムステップで
非圧縮条件を計算機精度水準で維持することを確認する．

%% TODO: exp11_3 の div(u) 時間推移をここに挿入
